\documentclass[11pt]{letter}
\usepackage[margin=1in]{geometry}
\usepackage{amsmath}
\usepackage{hyperref}

\signature{Ian Todd\\Sydney Medical School\\University of Sydney}
\address{Ian Todd\\Sydney Medical School\\University of Sydney\\Sydney, NSW 2006, Australia\\itod2305@uni.sydney.edu.au}

\begin{document}

\begin{letter}{The Editors\\Physical Review E\\American Physical Society}

\opening{Dear Editors,}

I am pleased to submit the manuscript ``Thermodynamic costs of dimensional reduction in stochastic dynamics'' for consideration as a Regular Article in \textit{Physical Review E}.

Landauer's principle provides the fundamental thermodynamic lower bound for logical information processing, linking bit erasure to heat dissipation. However, a wide class of physical and biological systems process information not by flipping discrete bits, but by projecting high-dimensional stochastic trajectories onto lower-dimensional continuous manifolds. In this work, we derive a thermodynamic bound for this process of dimensional reduction.

Using the framework of stochastic thermodynamics and overdamped Langevin dynamics, we identify a ``Dimensional Landauer Bound'':
\[
W_{\text{min}} \ge k_B T \ln 2 \, \Delta I + k_B T \, C_\Phi
\]
where $C_\Phi$ represents a geometric contraction cost determined by the Jacobian of the projection map and the curvature of the target manifold.

Our derivation shows that dimensional reduction is a geometrically irreversible operation that imposes an entropic cost distinct from logical erasure. We demonstrate the validity of this bound through analytical derivation and numerical simulations of Brownian motion on curved manifolds, coupled Kuramoto oscillators, and neural networks. We find that the maintenance of low-dimensional representations in non-equilibrium steady states requires continuous power dissipation proportional to the geometric complexity of the projection.

This work fits squarely within the scope of \textit{Physical Review E}, specifically the sections on \textbf{Statistical Physics} and \textbf{Biological Physics}. It extends established non-equilibrium work theorems (such as the Jarzynski equality) to the domain of manifold learning and neural dynamics, offering a physical mechanism for the efficiency of coherent biological states.

Given the interdisciplinary nature of the work, involving stochastic thermodynamics and information geometry, I suggest the following experts as potential reviewers:
\begin{enumerate}
\item \textbf{Udo Seifert} (University of Stuttgart) -- For expertise in stochastic thermodynamics.
\item \textbf{William Bialek} (Princeton University) -- For expertise in the physics of biological information.
\item \textbf{Christopher Jarzynski} (University of Maryland) -- For expertise in non-equilibrium work relations.
\item \textbf{Karl Friston} (University College London) -- For expertise in variational principles in neuroscience.
\end{enumerate}

I confirm that this manuscript has not been published and is not under consideration elsewhere. I have no conflicts of interest to declare.

\closing{Sincerely,}

\end{letter}
\end{document}
